%% ============================================================
%%  Chapter 5: Conclusion
%% ============================================================
\chapter{Conclusion}
\label{chap:conclusion}

% GUIDANCE: 4--6 pages. Write this last.
% Structure: summary → contributions revisited → limitations →
%            future work.
% Be honest about limitations -- committees respect this far more
% than hand-waving.


% ── 5.1 ─────────────────────────────────────────────────────────
\section{Summary}
\label{sec:conclusion_summary}

% Briefly restate what you set out to do and what you found.
% Do NOT introduce new results here.

\todo{Write 1--2 paragraphs summarising the thesis.
  Open with the problem (multi-robot coordination in warehouse
  environments under communication constraints).
  Summarise the approach (hybrid GCBBA + prioritised A*,
  event-driven dynamic replanning).
  Summarise the key findings without repeating Chapter 4 in full.}


% ── 5.2 ─────────────────────────────────────────────────────────
\section{Contributions Revisited}
\label{sec:conclusion_contributions}

% Map each contribution listed in Section~\ref{sec:intro_contributions}
% to the specific evidence from Chapter 4 that supports it.

\todo{Revisit each contribution from the Introduction and state
  which experimental result supports it. Be precise:
  ``Dynamic GCBBA achieves 100\% task completion at $\comm = 13$
  where Static GCBBA achieves 0\% (Section~\ref{subsec:exp_completion_rate})''.
  Avoid vague language like ``our method outperforms baselines''.}


% ── 5.3 ─────────────────────────────────────────────────────────
\section{Limitations}
\label{sec:conclusion_limitations}

% Be honest. Committees will ask about limitations regardless.
% Naming them yourself shows rigour.

\todo{Discuss limitations honestly. Suggested:
  (1) Centralised collision avoidance does not scale to
      truly decentralised settings (single point of failure
      in the path planner);
  (2) Experiments are simulation-only on a 30$\times$30 grid;
      real warehouses are larger and have different obstacle layouts;
  (3) SGA on disconnected graphs runs per component -- the
      sub-optimality ratio comparison is generous to SGA;
  (4) Fixed robot speeds and no sensor noise -- real robots
      have actuation uncertainty;
  (5) The 800-timestep ceiling may not be long enough for
      very high task loads at low comm range.}


% ── 5.4 ─────────────────────────────────────────────────────────
\section{Future Work}
\label{sec:conclusion_future}

\todo{Describe concrete future directions. Suggested:
  (1) Fully decentralised collision avoidance
      (distributed reservation tables or CBS-based);
  (2) Extension to time-windowed tasks (task deadlines);
  (3) Scaling to larger grids and higher agent counts;
  (4) Heterogeneous agents (different speeds, capacities);
  (5) Real robot implementation and evaluation;
  (6) Learning-based communication graph prediction to
      anticipate connectivity changes before they happen.}
